\documentclass[11pt,a4paper]{article}

% ---------------------------------------------------------
% PACKAGES
% ---------------------------------------------------------
\usepackage[utf8]{inputenc}
\usepackage[T1]{fontenc}
\usepackage{lmodern}
\usepackage{amsmath,amssymb,amsfonts,bm}
\usepackage{graphicx}
\usepackage{hyperref}
\usepackage{authblk}
\usepackage{geometry}
\usepackage{cite}
\usepackage{physics}
\usepackage{tensor}
\usepackage{enumitem}

\geometry{margin=1in}

% ---------------------------------------------------------
% TITLE
% ---------------------------------------------------------
\title{\textbf{Companion XV (v49.1)}\\[0.2cm]
\textbf{The Early-Universe Gravitational-Wave Spectrum}\\[0.15cm]
\large Relaxation-Induced Tensor Damping, Freeze-Out Geometry,
and CPT-Sector GW Signatures}

\author[1]{Michael Lehmann}
\affil[1]{\small Independent Researcher, Relaxation-Driven Cyclic Cosmology (RDCC)\\
\texttt{mi.lehmann@gmx.de}}

\date{September 2026}

\begin{document}
\maketitle

\begin{abstract}
The early-Universe gravitational-wave (GW) spectrum provides a uniquely
sensitive probe of the microphysical processes that shaped the primordial
plasma. In the Relaxation-Driven Cyclic Cosmology (RDCC), the relaxation
rate introduces an additional damping term in the tensor-mode equation,
suppressing primordial gravitational waves on scales that re-entered the
horizon during the relaxation-dominated epoch. In previous versions this
transition scale was written as \(k_{\mathrm{rel}}\sim a_{\mathrm{rel}}
H_{\mathrm{rel}}\). In the present Phase-1 integrated version the same
scale is inherited from the freeze-out geometry of Companion~XII,
\(k_{\mathrm{rel}}=a_{\mathrm{fr}}H(a_{\mathrm{fr}})\). Combined with a
possible sectoral temperature asymmetry and the logarithmic CPT modulation
of Companion~II, this yields a characteristic suppression feature in the
stochastic gravitational-wave background (SGWB). This Companion derives the
relaxation-corrected tensor evolution, the transfer function, and the early
Universe GW spectrum, and synchronizes the tensor sector with Companions~II,
XI, XII, VII, XLIV and~L.
\end{abstract}

\tableofcontents
\newpage

% ========================================================
\section{Motivation}
% ========================================================

The stochastic gravitational-wave background (SGWB) generated in the early
Universe provides a uniquely sensitive probe of high-energy physics. In
standard cosmology the SGWB is determined by the primordial tensor spectrum
and its subsequent evolution through radiation and matter domination. The
resulting spectrum is smooth, nearly scale-invariant, and free of dynamical
features except for small steps associated with changes in the relativistic
degrees of freedom.

In RDCC the situation is different. The relaxation mechanism modifies the
evolution of tensor modes during the earliest stages of cosmic history.
When the relaxation rate satisfies
\begin{equation}
  \Gamma(a)
  =
  \alpha(a)\, H(a)
  \gtrsim
  H(a),
\end{equation}
tensor modes experience an additional damping term that suppresses their
amplitude before or at horizon re-entry. This effect leaves a characteristic
imprint on the SGWB: a smooth, power-law suppression at frequencies
corresponding to modes that re-entered during the relaxation-dominated
epoch.

A second ingredient is a possible sectoral temperature asymmetry between the
two CPT-conjugate sectors. Different relaxation histories then imply
different damping scales. The precise temperature ratio is treated as an
open modeling assumption rather than a fixed input.

This Companion develops the early-Universe gravitational-wave sector of
RDCC and attaches it to the freeze-out geometry of Companion~XII. The SGWB
emerges as a central observational probe of the same transition that
governs growth suppression and gravitational slip.

The structure is as follows. Section~2 reviews tensor modes in standard
cosmology. Section~3 derives the relaxation-induced tensor damping.
Section~4 constructs the relaxation-corrected transfer function, including
the geometric identification of \(k_{\mathrm{rel}}\) and the logarithmic
CPT modulation. Section~5 presents the full early-Universe GW spectrum.
Section~6 develops observational predictions, and Section~7 concludes.

% ========================================================
\section{Connection to the Freeze-Out Geometry}
% ========================================================

Companion XII established the geometric origin of the relaxation scale
through the freeze-out of the global relaxation sector. The relaxation
scale is defined by
\begin{equation}
  k_{\mathrm{rel}}
  =
  a_{\mathrm{fr}} H(a_{\mathrm{fr}}),
\end{equation}
where
\begin{equation}
  a_{\mathrm{fr}}
  \simeq
  10^{-4}
\end{equation}
marks the freeze-out epoch of the relaxation field.

The tensor sector inherits this scale directly through horizon crossing
during the relaxation-dominated epoch. The resulting causal chain is
\begin{equation}
  \alpha
  \;\rightarrow\;
  a_{\mathrm{fr}}
  \;\rightarrow\;
  k_{\mathrm{rel}}
  \;\rightarrow\;
  T_{\mathrm{rel}}(k)
  \;\rightarrow\;
  \Omega_{\mathrm{GW}}(k).
\end{equation}
The stochastic gravitational-wave background therefore probes the same
freeze-out geometry that governs growth suppression and gravitational slip.

% ========================================================
\section{Tensor Modes in Standard Cosmology}
% ========================================================

In standard cosmology, primordial tensor perturbations obey the equation
\begin{equation}
  \ddot{h}_{k} + 3H \dot{h}_{k} + \frac{k^{2}}{a^{2}} h_{k} = 0,
\end{equation}
with two distinct dynamical regimes:
\begin{itemize}
  \item \textbf{Super-horizon regime} (\(k < aH\)):
        tensor modes freeze out with constant amplitude,
        \(\dot{h}_{k}\approx 0\).
  \item \textbf{Sub-horizon regime} (\(k > aH\)):
        tensor modes oscillate with physical frequency \(k/a\) and decay
        as \(a^{-1}\) due to Hubble friction.
\end{itemize}

The primordial tensor spectrum is commonly parameterized as
\begin{equation}
  \mathcal{P}_{h}(k)
  =
  A_{t}
  \left(
    \frac{k}{k_{*}}
  \right)^{n_{t}},
\end{equation}
where \(A_{t}\) is the tensor amplitude, \(n_{t}\) is the tensor tilt, and
\(k_{*}\) is a pivot scale.

The present-day gravitational-wave energy density per logarithmic interval
in wavenumber is
\begin{equation}
  \Omega_{\mathrm{GW}}(k)
  =
  \frac{1}{12}
  \left(
    \frac{k}{a_{0} H_{0}}
  \right)^{2}
  \mathcal{P}_{h}(k)\, T^{2}(k),
\end{equation}
where \(T(k)\) is the tensor transfer function.

In standard cosmology the transfer function is determined by horizon
re-entry during radiation or matter domination, changes in the relativistic
degrees of freedom, and the radiation--matter transition. Crucially,
standard cosmology contains no additional damping beyond Hubble friction.
The SGWB is therefore smooth and featureless, making it exceptionally
sensitive to any new early-Universe physics.

% ========================================================
\section{Relaxation-Induced Tensor Damping}
% ========================================================

In RDCC the evolution of tensor perturbations is modified by the
relaxation mechanism. The relaxation rate
\begin{equation}
  \Gamma(a)
  =
  \alpha(a)\, H(a)
\end{equation}
acts as an additional friction term in the tensor-mode equation. The
tensor-mode equation becomes
\begin{equation}
  \ddot{h}_{k}
  +
  \left( 3H + \Gamma \right) \dot{h}_{k}
  +
  \frac{k^{2}}{a^{2}} h_{k}
  =
  0,
\end{equation}
which can be written as
\begin{equation}
  \ddot{h}_{k}
  +
  (3 + \alpha) H\, \dot{h}_{k}
  +
  \frac{k^{2}}{a^{2}} h_{k}
  =
  0.
\end{equation}

\subsection{Super-Horizon Evolution}

For modes outside the horizon (\(k < aH\)), the gradient term is negligible
and the equation reduces to
\begin{equation}
  \ddot{h}_{k} + (3 + \alpha) H\, \dot{h}_{k} \approx 0.
\end{equation}
In standard cosmology tensor modes freeze out with constant amplitude. In
RDCC the additional relaxation term suppresses the amplitude according to
\begin{equation}
  h_{k}(a) \propto a^{-\alpha/2},
\end{equation}
implying that modes which remain super-horizon during the
relaxation-dominated epoch experience a smooth, power-law suppression.

\subsection{Sub-Horizon Evolution}

Once a mode re-enters the horizon (\(k > aH\)), the oscillatory term
dominates and the solution takes the form
\begin{equation}
  h_{k}(a)
  \propto
  \frac{1}{a}
  \exp\!
  \left[
    -\frac{1}{2}
    \int \alpha(a)\, d\ln a
  \right]
  \cos\!
  \left(
    k \int \frac{da}{a^{2} H}
  \right).
\end{equation}
The usual \(a^{-1}\) decay from Hubble friction is multiplied by an
additional exponential suppression factor sourced by the relaxation
dynamics.

\subsection{Relaxation-Dominated Epoch}

The relaxation-dominated regime is defined by
\begin{equation}
  \Gamma(a) \gtrsim H(a),
\end{equation}
or equivalently by a relaxation coefficient of order unity. The
relaxation-dominated epoch terminates at the freeze-out scale factor
\begin{equation}
  a_{\mathrm{rel}}
  \equiv
  a_{\mathrm{fr}}.
\end{equation}
Consequently, the relaxation scale becomes
\begin{equation}
  k_{\mathrm{rel}}
  =
  a_{\mathrm{fr}} H(a_{\mathrm{fr}}),
\end{equation}
in agreement with the freeze-out geometry established in Companion~XII.

Modes with
\begin{equation}
  k > k_{\mathrm{rel}}
\end{equation}
re-enter the horizon during the relaxation-dominated epoch and therefore
experience significant damping, while modes with
\begin{equation}
  k < k_{\mathrm{rel}}
\end{equation}
retain their standard amplitude. This produces a smooth, power-law
suppression feature in the SGWB.

\subsection{Sectoral Tensor Asymmetry}

A sectoral temperature asymmetry implies different relaxation histories in
the two sectors,
\begin{equation}
  \alpha_{(-)}(a) \neq \alpha_{(+)}(a).
\end{equation}
The precise sectoral temperature ratio remains an open modeling assumption.
Illustrative benchmark scenarios suggest modest sectoral temperature
asymmetries, while the qualitative prediction of sector-dependent tensor
histories remains unchanged. The sectoral asymmetry should therefore be
regarded as a testable feature rather than a fixed input parameter.

As a result:
\begin{itemize}
  \item tensor modes in the visible sector may experience stronger damping,
  \item tensor modes in the CPT-conjugate sector may experience weaker
        damping,
  \item the SGWB may exhibit a sectoral asymmetry in amplitude and spectral
        shape.
\end{itemize}

% ========================================================
\section{Relaxation-Corrected Transfer Function}
% ========================================================

The transfer function encodes the evolution of tensor modes from horizon
re-entry to the present day. In RDCC the relaxation mechanism introduces an
additional suppression factor at high frequencies. The
relaxation-corrected transfer function can be written as
\begin{equation}
  T^{2}(k)
  =
  T_{\mathrm{std}}^{2}(k)\,
  T_{\mathrm{rel}}^{2}(k),
\end{equation}
where \(T_{\mathrm{std}}(k)\) is the standard transfer function and
\(T_{\mathrm{rel}}(k)\) captures the relaxation-induced suppression.

\subsection{Standard Transfer Function}

The standard transfer function is given by
\begin{equation}
  T_{\mathrm{std}}^{2}(k)
  =
  \Omega_{r,0}
  \left[
    1
    +
    \frac{4}{3}
    \left(
      \frac{k}{k_{\mathrm{eq}}}
    \right)
    +
    \frac{5}{2}
    \left(
      \frac{k}{k_{\mathrm{eq}}}
    \right)^{2}
  \right]^{-1},
\end{equation}
where \(k_{\mathrm{eq}}\) is the wavenumber corresponding to
matter--radiation equality and \(\Omega_{r,0}\) is the present-day
radiation density fraction.

\subsection{Relaxation-Induced Suppression Factor}

The relaxation-induced suppression factor arises from the additional
damping term in the tensor-mode equation. For a mode with wavenumber
\(k\),
\begin{equation}
  T_{\mathrm{rel}}(k)
  =
  \exp\!
  \left[
    -\frac{1}{2}
    \int_{\ln a_{\mathrm{hc}}}^{\ln a_{\mathrm{fr}}}
    \alpha(a)\, d\ln a
  \right],
\end{equation}
where \(a_{\mathrm{hc}}\) is the scale factor at horizon crossing and
\(a_{\mathrm{fr}}\) marks the end of the relaxation-dominated epoch.
Modes that re-enter during the relaxation-dominated epoch experience the
largest suppression. Modes that re-enter later are unaffected.

\subsection{Characteristic Relaxation Scale}

The characteristic relaxation scale is
\begin{equation}
  k_{\mathrm{rel}}
  =
  a_{\mathrm{fr}} H(a_{\mathrm{fr}}).
\end{equation}
Modes with \(k > k_{\mathrm{rel}}\) re-enter during the
relaxation-dominated epoch and are therefore suppressed, while modes with
\(k < k_{\mathrm{rel}}\) retain their standard amplitude. The
relaxation-induced suppression factor can be approximated as
\begin{equation}
  T_{\mathrm{rel}}(k)
  \simeq
  \left(
    \frac{k}{k_{\mathrm{rel}}}
  \right)^{-\alpha_{\mathrm{eff}}},
  \qquad
  k > k_{\mathrm{rel}},
\end{equation}
where \(\alpha_{\mathrm{eff}}\) is an effective relaxation coefficient
that captures the integrated relaxation history.

\subsection{Geometric Interpretation of the Relaxation Scale}

The characteristic scale appearing in the tensor transfer function is no
longer treated as a phenomenological parameter. Within the modern RDCC
framework,
\begin{equation}
  k_{\mathrm{rel}}
  =
  a_{\mathrm{fr}} H(a_{\mathrm{fr}})
\end{equation}
is fixed by the freeze-out geometry of the relaxation field. Tensor
suppression therefore begins at the scale defined by the physical
transition between the relaxation-dominated and post-relaxation epochs.
The transfer function \(T_{\mathrm{rel}}(k)\) inherits its characteristic
scale from the same geometric transition that governs large-scale
structure observables.

\subsection{Sectoral Transfer Functions}

A possible sectoral temperature asymmetry implies different relaxation
histories in the two sectors and therefore sector-dependent transfer
functions,
\begin{equation}
  T_{\mathrm{rel}}^{(-)}(k)
  \neq
  T_{\mathrm{rel}}^{(+)}(k).
\end{equation}
The visible sector may exhibit a stronger high-frequency suppression,
while the CPT-conjugate sector may retain a larger tensor amplitude. The
precise ratio remains a modeling assumption, not a fixed input.

\subsection{CPT-Induced Logarithmic Modulation}

Companion~II demonstrated that tensor modes acquire a logarithmic phase
modulation due to the CPT structure of the bipartite Universe. The phase
shift is
\begin{equation}
  \Delta\phi(k)
  =
  \alpha_{\mathrm{IR}}
  \ln
  \left(
    \frac{k}{k_{b}}
  \right),
\end{equation}
leading to a tensor spectrum
\begin{equation}
  h(k)
  =
  h_{0}(k)
  \Biggl[
    1
    +
    \varepsilon(k)
    \cos
    \Bigl(
      \alpha_{\mathrm{IR}}
      \ln(k/k_{b})
      +
      \varphi_{\mathrm{CPT}}
    \Bigr)
  \Biggr].
\end{equation}
The complete RDCC tensor signal is therefore determined by both
\begin{itemize}
  \item relaxation-induced damping,
  \item CPT-induced logarithmic modulation.
\end{itemize}
The relaxation scale and the bounce scale represent two distinct physical
scales:
\begin{equation}
  k_{b} \neq k_{\mathrm{rel}}.
\end{equation}
The bounce scale controls the modulation structure, whereas the relaxation
scale controls the damping transition.

% ========================================================
\section{Full Early-Universe GW Spectrum}
% ========================================================

The present-day GW energy density per logarithmic interval in wavenumber
is
\begin{equation}
  \Omega_{\mathrm{GW}}(k)
  =
  \frac{1}{12}
  \left(
    \frac{k}{a_{0} H_{0}}
  \right)^{2}
  \mathcal{P}_{h}(k)\,
  T_{\mathrm{std}}^{2}(k)\,
  T_{\mathrm{rel}}^{2}(k).
\end{equation}
The resulting spectrum exhibits three characteristic regimes.

\subsection{Low-Frequency Regime}

Modes with \(k < k_{\mathrm{rel}}\) re-enter the horizon after the
relaxation-dominated epoch. Their amplitude is unaffected by the
relaxation mechanism, and the spectrum reduces to the standard result:
\begin{equation}
  \Omega_{\mathrm{GW}}(k)
  \simeq
  \Omega_{\mathrm{GW}}^{\mathrm{std}}(k),
  \qquad
  k < k_{\mathrm{rel}}.
\end{equation}

\subsection{Intermediate Regime}

Modes with \(k \sim k_{\mathrm{rel}}\) re-enter near freeze-out. These
modes experience partial suppression, leading to a smooth transition. The
spectrum takes the form
\begin{equation}
  \Omega_{\mathrm{GW}}(k)
  \simeq
  \Omega_{\mathrm{GW}}^{\mathrm{std}}(k)
  \left(
    \frac{k}{k_{\mathrm{rel}}}
  \right)^{-2\alpha_{\mathrm{eff}}},
  \qquad
  k \sim k_{\mathrm{rel}}.
\end{equation}

\subsection{High-Frequency Regime}

Modes with \(k > k_{\mathrm{rel}}\) re-enter during the
relaxation-dominated epoch and experience the strongest suppression:
\begin{equation}
  \Omega_{\mathrm{GW}}(k)
  \simeq
  \Omega_{\mathrm{GW}}^{\mathrm{std}}(k)
  \left(
    \frac{k}{k_{\mathrm{rel}}}
  \right)^{-2\alpha_{\mathrm{eff}}},
  \qquad
  k > k_{\mathrm{rel}}.
\end{equation}
This power-law suppression is the defining signature of
relaxation-induced tensor damping. Superimposed on it, the logarithmic
modulation of Companion~II may appear as an oscillatory envelope.

\subsection{Sectoral GW Spectra}

Different relaxation histories in the two sectors imply
\begin{equation}
  \Omega_{\mathrm{GW}}^{(+)}(k)
  \neq
  \Omega_{\mathrm{GW}}^{(-)}(k).
\end{equation}
The visible sector may exhibit a stronger high-frequency suppression,
while the CPT-conjugate sector may retain a larger tensor amplitude. The
qualitative prediction of sector-dependent spectra is robust; the
numerical temperature ratio is not fixed in v49.1.

% ========================================================
\section{Observational Predictions}
% ========================================================

The relaxation-induced tensor damping and a possible sectoral temperature
asymmetry lead to distinctive signatures in the SGWB.

\subsection{High-Frequency Suppression}

The defining prediction is a smooth, power-law suppression of the SGWB at
high frequencies:
\begin{equation}
  \Omega_{\mathrm{GW}}(k)
  \propto
  \left(
    \frac{k}{k_{\mathrm{rel}}}
  \right)^{-2\alpha_{\mathrm{eff}}},
  \qquad
  k > k_{\mathrm{rel}},
\end{equation}
with
\begin{equation}
  k_{\mathrm{rel}}
  =
  a_{\mathrm{fr}} H(a_{\mathrm{fr}}).
\end{equation}

\subsection{Sectoral GW Asymmetry}

A sectoral temperature asymmetry, if present, implies different
relaxation histories. The visible sector may exhibit a stronger
high-frequency suppression, while the CPT-conjugate sector may retain a
larger tensor amplitude. This asymmetry is a testable feature, not a
fixed input.

\subsection{Connection to Structure-Formation Observables}

The same freeze-out scale appears in the growth factor and growth index
of Companions~VII and~L and in the gravitational slip of Companion~XLIV.
The SGWB therefore provides a complementary probe of one and the same
geometric transition.

\subsection{Experimental Prospects}

Future gravitational-wave observatories will be able to probe the
relaxation-induced features of the SGWB:
\begin{itemize}
  \item \textbf{LISA}: sensitive to the intermediate regime near
        \(k_{\mathrm{rel}}\),
  \item \textbf{DECIGO/BBO}: sensitive to the high-frequency suppression,
  \item \textbf{Einstein Telescope / Cosmic Explorer}: complementary
        constraints at lower frequencies.
\end{itemize}
A detection of a smooth, power-law suppression at high frequencies,
possibly with a logarithmic modulation, would provide strong evidence for
relaxation-driven tensor damping.

% ========================================================
\section{Conclusion}
% ========================================================

In this Companion we have developed the early-Universe gravitational-wave
sector of RDCC. The relaxation mechanism modifies the evolution of tensor
perturbations and introduces an additional damping term on scales that
re-entered the horizon during the relaxation-dominated epoch.

The modernized tensor sector is fully compatible with the freeze-out
geometry developed in Companion~XII. The characteristic relaxation scale
is no longer inserted phenomenologically but identified through
\begin{equation}
  k_{\mathrm{rel}}
  =
  a_{\mathrm{fr}} H(a_{\mathrm{fr}}).
\end{equation}
The resulting tensor hierarchy becomes
\begin{equation}
  \alpha
  \;\rightarrow\;
  a_{\mathrm{fr}}
  \;\rightarrow\;
  k_{\mathrm{rel}}
  \;\rightarrow\;
  T_{\mathrm{rel}}(k)
  \;\rightarrow\;
  \Omega_{\mathrm{GW}}(k).
\end{equation}

The SGWB therefore exhibits a smooth, power-law suppression at high
frequencies, with the transition fixed by freeze-out. A possible sectoral
temperature asymmetry leads to sector-dependent damping and is treated as
a testable modeling assumption. Logarithmic CPT modulation, controlled by
the distinct bounce scale \(k_{b}\), may overlay the damping envelope.

Together with Companions~II, XI, XII, VII, XLIV and~L, these results
establish the SGWB as a central observational probe of the same
relaxation-driven geometry that underlies growth and gravitational slip.
Future observatories such as LISA, DECIGO, BBO and the Einstein Telescope
can test the suppression feature and, if present, the sectoral asymmetry
and logarithmic modulation.

\end{document}